\newpage\section{坐标}

本节概述用户可用的各种坐标系。既然缩放操作已从 \tkzname{Lua} 层移除，这样的说明就显得比以往任何时候都更有必要。

\subsection{坐标的一般用法}
与 \tkzNamePack{tkz-euclide} 一样，\tkzNamePack{tkz-elements} 建立在二维正交笛卡尔坐标系之上，默认单位为厘米。

在一个专注于欧几里得几何的语境中使用我们将要介绍的 \emph{occs}（正交坐标系）会是不协调的。

此外，在 \tkzNamePack{tkz-elements} 中，点的概念与一个复数的附标绑定在一起。为保持代码清晰一致，在 Lua 内部修改单位的选项已被刻意省略。

\medskip

\textbf{结论：} 对于用 \tkzNamePack{tkz-elements} 创建的任何图形，所有点都放在一个以厘米为单位的二维正交坐标系中。

   \begin{tikzpicture}
   \pgfkeys{/pgf/number format/.cd,std,precision=2}
   \let\pmpn\pgfmathprintnumber
   \tkzDefPoints{2/3/M,0/0/O,2/0/A,0/3/B}
   \tkzLabelPoints(O)
   \tkzLabelPoint[right](M){$M: z_M = 2 + 3i$}
   \tkzDrawPoints(M,O)
   \tkzPointShowCoord(M)
   \tkzLabelPoint[below,teal](A){$2$}
   \tkzLabelPoint[left,teal](B){$3$}
   \tkzDrawSegments[->,add = 0 and 0.25](O,B O,A)
   \begin{scope}[every annotation/.style={fill=lightgray!15,anchor = east}]
   \node [annotation,font =\small,text width=6cm] at (current bounding box.west) {
$M$ 的坐标
   \begin{mybox}{}
   \code{ z.A = point(2, 3)}
   \end{mybox}
       };
   \end{scope}
   \end{tikzpicture}

\subsection{重心坐标的使用}

重心坐标描述一个点相对于参考三角形的位置。平面上任意一点都可以用重心坐标表示，这些坐标在相差一个标量倍数（位似）的意义下确定。或者，也可以把它们归一化，使三者之和等于 1。

重心坐标在三角形几何中特别有用，尤其是在分析那些在仿射变换下不变的性质——即不依赖于角度的性质——时。

考虑三角形 $ABC$。人们可以用重心坐标来定义重心、垂心等关键点。

\medskip
\texttt{关于垂心。}

计算垂心的重心坐标需要知道三角形的三个角。它们作为属性存放在表 \texttt{T.ABC} 中：
\begin{itemize}
  \item \texttt{T.ABC.alpha} —— 顶点 $A$ 处的角 $\widehat{A}$
  \item \texttt{T.ABC.beta}  —— 顶点 $B$ 处的角 $\widehat{B}$
  \item \texttt{T.ABC.gamma} —— 顶点 $C$ 处的角 $\widehat{C}$
\end{itemize}

见 [\ref{ssub:_triangle_barycentric_ka_kb_kc}]

\medskip
\texttt{取回重心坐标。}

可以计算一个给定点相对于某个三角形的重心坐标。返回的值会自动归一化。用法见第 [\ref{ssub:_triangle_barycentric__coordinates_pt}] 节。


\begin{minipage}{.5\textwidth}
\directlua{
z.A = point(0, 0)
z.B = point(3, 1)
z.C = point(1, 3)
T.ABC = triangle(z.A, z.B, z.C)
z.G = T.ABC:barycentric(1,1,1)
z.H = T.ABC:barycentric(math.tan(T.ABC.alpha),
                        math.tan(T.ABC.beta),
                        math.tan(T.ABC.gamma))}
\begin{tikzpicture}
\tkzGetNodes
\tkzDrawPolygon(A,B,C)
\tkzDrawPoints(A,B,C,G,H)
\tkzLabelPoints(A,B,G)
\tkzLabelPoints[above](C,H)
\end{tikzpicture}
\end{minipage}
\begin{minipage}{.5\textwidth}
\begin{tkzexample}[code only]
\directlua{
z.A = point(0, 0)
z.B = point(3, 1)
z.C = point(1, 3)
T.ABC = triangle(z.A, z.B, z.C)
z.G = T.ABC:barycentric(1, 1, 1)
z.H = T.ABC:barycentric(math.tan(T.ABC.alpha),
                        math.tan(T.ABC.beta),
                        math.tan(T.ABC.gamma))}
\end{tkzexample}
\end{minipage}

\subsection{三线坐标的使用}
\label{sub:use_of_trilinear_coordinates}

点 $P$ 相对于三角形 $ABC$ 的三线坐标，是一组与 $P$ 到三角形各边的有向距离成比例的三个数。这些坐标是齐次的，通常写作 $x : y : z$ 或 $(x, y, z)$。

由于只有坐标之间的比值才有意义，三线坐标特别适合表达那些在缩放不变的几何关系。

于是在下一个示例中有 $ a' : b' : c' = ka : kb : kc$（其中 $a = BC, b = AC, c = AB$）。

\begin{minipage}{.5\textwidth}
\directlua{
z.A = point(0, 0)
z.B = point(5, 0)
z.C = point(1, 3)
T.ABC = triangle(z.A, z.B, z.C)
z.L = T.ABC:trilinear(T.ABC.a,
                      T.ABC.b,
                      T.ABC.c)
z.a, z.b, z.c = T.ABC:projections(z.L)}
\begin{tikzpicture}
\tkzGetNodes
\tkzDrawPolygon(A,B,C)
\tkzDrawSegments(L,a L,b L,c)
\tkzDrawPoints(A,B,C,L)
\tkzLabelPoints(A,B)
\tkzLabelPoints[above](C,L)
\tkzLabelSegment[left](L,a){$ka$}
\tkzLabelSegment[above](L,b){$kb$}
\tkzLabelSegment[right](L,c){$kc$}
\end{tikzpicture}
\end{minipage}
\begin{minipage}{.5\textwidth}
\begin{tkzexample}[code only]
\directlua{
z.A = point(0, 0)
z.B = point(5, 0)
z.C = point(1, 3)
T.ABC = triangle(z.A, z.B, z.C)
z.L = T.ABC:trilinear(T.ABC.a,
                      T.ABC.b,
                      T.ABC.c)
z.a, z.b, z.c = T.ABC:projections(z.L)}
\end{tkzexample}
\end{minipage}

\subsection{OCCS}

该类的对象就是 \emph{正交笛卡尔坐标系}。它们由参考坐标系经过平移和旋转得到。可用于简化某些表达式和坐标。
见 [\ref{sec:orthonormal_cartesian_coordinate_system}
]
\endinput